\scanpage{84}
\mspage{6}
% archive foliation: 83

\deleted{On définit \unsure{des}}

Le \deleted{catégorie}\inserted{champ} en question n'est pas rigide --- i.e. un objet peut avoir des \deleted{\unsure{autom}} non triviaux. \deleted{Les}\inserted{Un} autom. \deleted{sont}\inserted{d'un objt (34) est} défini\deleted{\unsure{s}} par des autom de $\underline{J}$, i.e. une section du schéma en groupes étale $\unsure{\mathfrak{S}|\mathcal G}_{\underline{J}}$, qui sont compatible avec les données b) et c). La compatibilité avec c) signifie que c'est une section de
\[
\unsure{\mathfrak{S}|\mathcal G}_{\underline{J}}^+\simeq\deleted{\unsure{\operatorname{id}}}\underbrace{S}_{\substack{\text{correspond}\\\text{à la section}\\\text{unité du}\\\text{schéma en groupes}}}\amalg\underline{\omega}_{\underline{J}}
\]
une section qui est $\ne\operatorname{id}$ sur \unsure{tte} fibre géom. revient donc à une section de $\underline{\omega}_{\underline{J}}$, soit $\rho$ --- de sorte que $\underline{J}$ devient un torseur sous $(\mathbb{Z}/3\mathbb{Z})_S$. La compatibilité avec b) signifie que $t=\rho(t)$ i.e. que $t$ \deleted{correspond à}\inserted{est une} section du sous-schéma $\underline{P}_{\underline{J}}\subset\Sigma_{\underline{J}}^*$,
\margin{$\underline{P}_{\underline{J}}=(\Sigma_{\underline{J}})^{\unsure{\mathfrak{S}|\mathcal G}_{\underline{J}}^+}$}
\deleted{\unsure{est}\illegible{3}}
